\documentclass[aspectratio=169]{beamer}
\usepackage[utf8]{inputenc}
\usepackage[spanish]{babel}
\usepackage{graphicx}
\usepackage{booktabs}
\usepackage{amsmath}
\usepackage{tikz}
\usepackage{xcolor}

\usetheme{Madrid}
\usecolortheme{default}

\title{Experimento de Ablación Sistemática}
\subtitle{Validación de Especificidad Clínica en Detección de Espasticidad Infantil}
\author{Máximo Fernández Riera}
\institute{Universitat Oberta de Catalunya (UOC)}
\date{Abril 2026}

\begin{document}

\frame{\titlepage}

\begin{frame}{Índice}
\tableofcontents
\end{frame}

\section{Concepto}

\begin{frame}{¿Qué es la Ablación?}
\begin{block}{Definición}
Técnica metodológica derivada de la neurociencia que consiste en \textbf{desactivar sistemáticamente componentes} de un sistema para determinar su contribución individual al rendimiento global.
\end{block}

\vspace{0.5cm}

\begin{columns}
\column{0.5\textwidth}
\textbf{Neurociencia:}
\begin{itemize}
\item Lesionar región cerebral
\item Observar cambio conductual
\item Inferir función
\end{itemize}

\column{0.5\textwidth}
\textbf{Machine Learning:}
\begin{itemize}
\item Eliminar componente del modelo
\item Medir cambio en rendimiento
\item Inferir importancia
\end{itemize}
\end{columns}
\end{frame}

\section{Contexto del Experimento}

\begin{frame}{Pipeline de Generación Sintética}
\begin{center}
\begin{tikzpicture}[node distance=2cm, auto]
\node[draw, rectangle, fill=blue!20] (video) {Vídeo Normal};
\node[draw, rectangle, fill=orange!20, right of=video, xshift=2cm] (pert) {7 Perturbaciones};
\node[draw, rectangle, fill=green!20, right of=pert, xshift=2cm] (patho) {Vídeo Patológico};

\draw[->, thick] (video) -- (pert);
\draw[->, thick] (pert) -- (patho);
\end{tikzpicture}
\end{center}

\vspace{0.5cm}

\begin{block}{7 Perturbaciones Clínicamente Motivadas}
\begin{enumerate}
\item \texttt{amplitude\_scale} $\rightarrow$ Rigidez muscular
\item \texttt{apply\_tremor} $\rightarrow$ Temblor patológico
\item \texttt{apply\_asymmetry} $\rightarrow$ Asimetría de movimiento
\item \texttt{apply\_motion\_blur} $\rightarrow$ Limitación de movimiento
\item \texttt{apply\_temporal\_jitter} $\rightarrow$ Fluctuación temporal
\item \texttt{apply\_global\_noise} $\rightarrow$ Ruido muscular
\item \texttt{apply\_common\_acquisition\_artifacts} $\rightarrow$ Defectos de captura
\end{enumerate}
\end{block}
\end{frame}

\section{Metodología}

\begin{frame}{Diseño Experimental}
\begin{block}{Condiciones del Experimento (8 total)}
\begin{itemize}
\item \textbf{A\_completo}: Todas las perturbaciones activas (baseline)
\item \textbf{B\_sin\_rigidez}: Excluye \texttt{amplitude\_scale}
\item \textbf{C\_sin\_temblor}: Excluye \texttt{apply\_tremor}
\item \textbf{D\_sin\_asimetria}: Excluye \texttt{apply\_asymmetry}
\item \textbf{E\_sin\_motion\_blur}: Excluye \texttt{apply\_motion\_blur}
\item \textbf{F\_sin\_jitter}: Excluye \texttt{apply\_temporal\_jitter}
\item \textbf{G\_sin\_ruido}: Excluye \texttt{apply\_global\_noise}
\item \textbf{H\_sin\_artifacts}: Excluye \texttt{apply\_common\_acquisition\_artifacts}
\end{itemize}
\end{block}
\end{frame}

\begin{frame}{Protocolo de Ejecución}
\begin{enumerate}
\item \textbf{Para cada condición:}
\begin{itemize}
\item Generar dataset sintético con perturbaciones específicas
\item Entrenar 4 modelos (LR, RF, SVM, XGBoost)
\item Evaluar en conjunto de test
\item Registrar métricas de rendimiento
\end{itemize}

\vspace{0.3cm}

\item \textbf{Análisis comparativo:}
\begin{itemize}
\item Calcular $\Delta$ accuracy = accuracy\textsubscript{baseline} - accuracy\textsubscript{ablated}
\item Calcular impacto relativo: $\frac{\Delta}{\text{accuracy}_{\text{baseline}}} \times 100$
\item Ranking de perturbaciones por impacto
\end{itemize}

\vspace{0.3cm}

\item \textbf{Visualización:}
\begin{itemize}
\item Tabla comparativa de métricas
\item Gráfico de barras de impacto
\item Tabla LaTeX para tesis
\end{itemize}
\end{enumerate}
\end{frame}

\section{Interpretación}

\begin{frame}{Análisis de Resultados}
\begin{block}{Interpretación del $\Delta$ Accuracy}
\begin{itemize}
\item \textcolor{green!70!black}{\textbf{$\Delta \approx 0$}}: La perturbación excluida \textbf{no aporta} información discriminativa significativa
\begin{itemize}
\item[$\rightarrow$] El modelo detecta otras manifestaciones clínicas
\item[$\rightarrow$] Robustez ante variabilidad
\end{itemize}

\vspace{0.3cm}

\item \textcolor{red!70!black}{\textbf{$\Delta > 0$ (grande)}}: La perturbación excluida \textbf{contribuye significativamente}
\begin{itemize}
\item[$\rightarrow$] Componente clave para discriminación
\item[$\rightarrow$] Posible dependencia excesiva
\end{itemize}
\end{itemize}
\end{block}

\vspace{0.3cm}

\begin{alertblock}{Escenario Ideal}
Todas las condiciones mantienen accuracy $> 0.75$ $\rightarrow$ El modelo ha aprendido un \textbf{patrón multidimensional robusto}, no un único artefacto.
\end{alertblock}
\end{frame}

\begin{frame}{Validación de Especificidad Clínica}
\begin{center}
\begin{tikzpicture}
\node[draw, rectangle, fill=red!20, text width=4cm, align=center] at (0,2) (bad) {\textbf{Modelo Débil}\\ Aprende artefactos\\ específicos};
\node[draw, rectangle, fill=green!20, text width=4cm, align=center] at (6,2) (good) {\textbf{Modelo Robusto}\\ Aprende manifestaciones\\ clínicas reales};

\node[text width=3.5cm, align=center] at (0,0) {$\Delta$ accuracy muy alto\\ al excluir 1-2 perturbaciones};
\node[text width=3.5cm, align=center] at (6,0) {$\Delta$ accuracy bajo\\ en todas las condiciones};

\draw[->, thick, red] (bad) -- (0,0.5);
\draw[->, thick, green!70!black] (good) -- (6,0.5);
\end{tikzpicture}
\end{center}
\end{frame}

\section{Objetivo en el TFM}

\begin{frame}{Contribución al Trabajo}
\begin{block}{Objetivo Principal}
Demostrar que el pipeline de generación sintética produce datos que capturan \textbf{características clínicas generalizables} de la espasticidad infantil, no solo parámetros artificiales de perturbación.
\end{block}

\vspace{0.5cm}

\begin{columns}
\column{0.5\textwidth}
\textbf{Control de Validez:}
\begin{itemize}
\item Validez interna
\item Validez de constructo
\item Especificidad clínica
\end{itemize}

\column{0.5\textwidth}
\textbf{Aportación Científica:}
\begin{itemize}
\item Metodología rigurosa
\item Transparencia experimental
\item Reproducibilidad
\end{itemize}
\end{columns}

\vspace{0.5cm}

\begin{alertblock}{Implicación}
Si el modelo mantiene rendimiento aceptable en todas las condiciones de ablación, se refuerza la \textbf{validez ecológica} del enfoque sintético para un contexto clínico real.
\end{alertblock}
\end{frame}

\begin{frame}{Ejemplo de Resultados Esperados}
\begin{table}
\centering
\small
\begin{tabular}{lccc}
\toprule
\textbf{Condición} & \textbf{Accuracy (SVM)} & \textbf{$\Delta$ (pp)} & \textbf{Impacto rel.} \\
\midrule
Completo & 0.8274 & - & 100\% \\
Sin rigidez & 0.8110 & -1.64 & 98.0\% \\
Sin temblor & 0.8208 & -0.66 & 99.2\% \\
Sin asimetría & 0.8046 & -2.28 & 97.2\% \\
Sin motion blur & 0.8175 & -0.99 & 98.8\% \\
Sin jitter & 0.8241 & -0.33 & 99.6\% \\
Sin ruido & 0.8143 & -1.31 & 98.4\% \\
Sin defectos captura & 0.8208 & -0.66 & 99.2\% \\
\bottomrule
\end{tabular}
\caption{Resultados hipotéticos del experimento de ablación.}
\end{table}

\vspace{0.3cm}
\textbf{Interpretación:} Todas las condiciones mantienen accuracy $> 0.80$ $\rightarrow$ Robustez confirmada.
\end{frame}

\section{Conclusiones}

\begin{frame}{Conclusiones}
\begin{enumerate}
\item La ablación sistemática es una técnica de \textbf{validación interna} esencial en ML aplicado a salud.

\vspace{0.3cm}

\item Permite distinguir entre:
\begin{itemize}
\item Modelos que aprenden \textcolor{red}{artefactos específicos} (overfitting a perturbaciones)
\item Modelos que aprenden \textcolor{green!70!black}{patrones clínicos robustos}
\end{itemize}

\vspace{0.3cm}

\item En nuestro TFM, demuestra que el pipeline sintético genera datos con \textbf{validez clínica}, no solo técnica.

\vspace{0.3cm}

\item Refuerza la \textbf{credibilidad científica} del trabajo ante evaluadores y comunidad médica.
\end{enumerate}
\end{frame}

\begin{frame}
\begin{center}
\Huge ¿Preguntas?
\end{center}
\end{frame}

\end{document}
